%----------------------------------------------------------------------
\section{The Discrimination Framework}\label{sec:framework}\footnote{Section-level Toulmin. \emph{Move}: establish the algebraic vocabulary. \emph{ZFC comparison}: ZFC takes membership as primitive and undefined; the framework takes discrimination (a binary predicate over $\GF(2)$) as primitive, giving membership algebraic content from the start. \emph{Residue instance}: the four-cell structure is the first appearance of gap/overlap --- the residue mechanism that recurs at every subsequent level.}
%----------------------------------------------------------------------

\subsection{Discriminators}

\begin{definition}[Discriminator]\label{def:discriminator}
A \emph{discriminator} is a grade-1 element $d \in \bigwedge^1 V$
together with a single predicate $d: S \to \GF(2)$. The
predicate \emph{fires} ($d(a) = 1$) or is \emph{silent}
($d(a) = 0$) on each element $a \in S$. A discriminator
is \emph{keyed to} a property of the data and tests whether
that property is active for a given datum. It provides one
bit of information per element.
\end{definition}

\begin{definition}[$\tau$ and $\sigma$ as datum properties]
\label{def:tau-sigma}
$\tau$ and $\sigma$ are two properties that a datum may
possess. They are attributes of the \emph{data}, not of
the discriminators. A discriminator keyed to $\tau$ tests
whether $\tau$ is active for a datum; a discriminator
keyed to $\sigma$ tests whether $\sigma$ is active.

For a datum $a$: if a $\tau$-keyed discriminator fires,
$a$ carries the $\tau$ label ($\tau(a) = 1$). If a
$\sigma$-keyed discriminator fires, $a$ carries the
$\sigma$ label ($\sigma(a) = 1$). The two labels are
independent --- a datum may carry both, one, or neither.
\end{definition}

\begin{definition}[Membership profile]\label{def:profile}
The \emph{membership profile} of a datum $a$ is the pair
\begin{equation}
  (\tau(a),\; \sigma(a)) \in \GF(2) \times \GF(2),
\end{equation}
recording which properties have been recognized as active
by the discriminators applied so far. Determining the
full profile requires (at least) two discriminators: one
keyed to $\tau$, one keyed to $\sigma$.

The profile is not merely a recorded tuple: it is itself
a \emph{composed discriminator}. The pair $(d_\tau, d_\sigma)$
--- a $\tau$-keyed discriminator and a $\sigma$-keyed
discriminator applied to the same datum --- corresponds
to the wedge product $d_\tau \wedge d_\sigma \in
\bigwedge^2 V$. The pair compares the results of two
base discriminators, and the comparison is the grade-2
element that the exterior algebra provides.

The $S_3$ symmetry of the triangle identity
(Definition \ref{def:triangle}) acts on these
compositions: rotating $\kappa \to \sigma \to \tau$
changes which discriminators are being compared
and in what role. The six adjunctions of Theorem
\ref{thm:adjunction} are six ways of composing
discriminator pairs.
\end{definition}

\begin{remark}[The four-cell structure]
The membership profile takes values in $\GF(2) \times \GF(2)$,
giving four possible states per datum:
\begin{center}
\begin{tabular}{c|c|l}
$\tau(a)$ & $\sigma(a)$ & Interpretation \\
\hline
$1$ & $0$ & $\tau$ active, $\sigma$ inactive: \emph{ordered} \\
$0$ & $1$ & $\sigma$ active, $\tau$ inactive: \emph{ordered} \\
$0$ & $0$ & both inactive: the \emph{gap} \\
$1$ & $1$ & both active: the \emph{overlap}
\end{tabular}
\end{center}
The four cells are the product of two independent
properties of the data. An element is \emph{discriminated}
(ordered) when exactly one property is active. It is
\emph{non-discriminated} when both are active (overlap)
or neither is (gap).

The overlap $(1,1)$ arises in two situations:
\begin{itemize}[nosep]
  \item \textbf{Genuine over-determination}: both
    properties legitimately hold for the datum, signaling
    that $\kappa$ needs to evolve (a finer discriminator
    is needed to separate what the current regime conflates).
  \item \textbf{Truncated projection}: an artifact of
    viewing a finer discrimination at coarser resolution.
\end{itemize}
Both produce the same algebraic signature ($\chi = 0$)
and both are residues.

A \emph{higher-level discriminator} can test composite
properties --- for instance, the XOR signature
$\chi(a) = \tau(a) + \sigma(a)$ tests whether exactly
one property is active. Such discriminators depend on
the results of base discriminators and are therefore
\emph{dependent type constructions}: the type of the
higher-level predicate depends on the values produced
by the base predicates. The Boolean dimension
(Definition \ref{def:boolean-dim}) is the principal
example: $\chi(a) = 1$ for all $a$ is a dependent type
that requires both base discriminators to have been
applied.

All four cells are \emph{post-interrogation} states: writing
$(\tau(a), \sigma(a)) = (0,0)$ means the relevant
discriminators have been applied and both properties were
found inactive. The value $0$ is a positivist assertion:
``I checked and found no match.'' Representing the
pre-interrogation state (no discriminator applied)
requires the $\GF(2)^3$ expansion with coverage bit
$\gamma$ (Section \ref{sec:expanded}).
\end{remark}

\begin{definition}[Profile linearity]\label{def:profile-linear}
The map from discriminators to their predicates is
\emph{linear}: for discriminators $d_1, d_2 \in
\bigwedge^1 V$ and any element $a \in S$,
\begin{equation}
  (d_1 + d_2)(a) = d_1(a) + d_2(a),
\end{equation}
where addition on the left is in $V$ and on the right is
in $\GF(2)$. This is a structural axiom: the predicate of
a sum of discriminators is the sum of the individual
predicates. It ensures that the evaluation map
$d \mapsto d(\cdot)$ respects the vector space structure
of $V$.
\end{definition}

\begin{remark}
Profile linearity (Definition \ref{def:profile-linear}) is not a free consequence of the
definitions --- it is an explicit requirement that ties
the algebraic structure of $V$ to the operational behavior
of discriminators. Without it, discriminators would be
elements of $V$ algebraically but their predicates could
behave non-linearly, severing the connection between the
exterior algebra and the discrimination process. Profile
linearity is what makes the algebra \emph{faithful} to
the operations it encodes.
\end{remark}

The semantic distinctions underlying these definitions
--- single predicate, datum properties, composed pairs,
dependent types --- are verified in
\S\ref{app:discriminator-semantics}.

\begin{definition}[XOR signature]
The \emph{XOR signature} of an element $a$ under a
discriminator pair $(d_\tau, d_\sigma)$ is
\begin{equation}
  \chi(a) = d_\tau(a) + d_\sigma(a) \in \GF(2).
\end{equation}
When $\chi(a) = 1$, exactly one predicate fires:
the element is \emph{discriminated} (ordered).
When $\chi(a) = 0$, both fire or neither fires:
the element is \emph{non-discriminated}.
\end{definition}

\begin{definition}[Residue]\label{def:residue}
The \emph{residue} of a discriminator pair
$(d_\tau, d_\sigma)$ applied to a collection $S$ is
\begin{equation}
  R(S) = \{ a \in S \mid \chi(a) = 0 \}.
\end{equation}
These are the elements for which discrimination was
non-trivial in $\GF(2)$: $1 + 1 = 0$ (overlap) or
$0 + 0 = 0$ (gap).
\end{definition}

\begin{remark}[Two origins of residue]\label{rem:residue-origins}
Gap residues ($\tau = \sigma = 0$) arise from \emph{novelty}: the
element has no membership under either predicate. Overlap residues
($\tau = \sigma = 1$) arise from \emph{abstraction}: the element
is above the discriminator's resolution, as when presenting ``mammal''
to a cat/dog discriminator. Both are zeros in $\GF(2)$, but they
demand $\kappa$-evolution in opposite directions: gaps extend the
tree laterally (new specificity); overlaps extend the tree upward
(new generality).
\end{remark}

\subsection{The discrimination regime $\kappa$}

\begin{definition}[Discrimination regime]\label{def:kappa}
A \emph{discrimination regime} $\kappa$ is a finite, ordered
sequence of discriminators $(d_1, d_2, \ldots, d_n)$, each
providing a single predicate $d_i: S \to \GF(2)$. The
\emph{dimensionality} of $\kappa$ is $n$. The $\tau/\sigma$
pairing structure (Definition \ref{def:profile}) is imposed
on pairs of discriminators within $\kappa$; it is not
intrinsic to individual discriminators.
\end{definition}

\begin{definition}[$\kappa$-evolution]\label{def:evolution}
A \emph{$\kappa$-evolution} is the extension of $\kappa$ by a new
discriminator $d_{n+1}$, producing $\kappa' = (d_1, \ldots, d_n, d_{n+1})$.
Evolution is triggered by the articulation of a predicate, which is
simultaneously the construction of a residue via the wedge product.
\end{definition}

\begin{proposition}[Monotonicity of $\kappa$]\label{prop:monotone}
$\kappa$-evolution is strictly monotonic: if $\kappa \to \kappa'$,
then $\dim(\kappa') = \dim(\kappa) + 1$. Dimensions are never removed.
\end{proposition}

\begin{proof}
By construction. Each evolution appends a new discriminator. The write
path is linear in the resource-consumption sense: the residue that
triggered the evolution is consumed and cannot trigger a second evolution.
\end{proof}

\subsection{Write and read regimes}\label{sec:write-read}

\begin{definition}[Write path]
The \emph{write path} is the $\kappa$-evolution process: each
residue is consumed exactly once to produce exactly one new
dimension.
\end{definition}

\begin{definition}[Read path]
The \emph{read path} is the process of evaluating discriminator
predicates over a collection $S$ given a fixed $\kappa$. Each
discriminator $d_i$ produces one bit $d_i(a) \in \GF(2)$ per
element. For a $\tau/\sigma$ pair $(d_\tau, d_\sigma)$ within
$\kappa$, the read path produces the membership profile
$(d_\tau(a), d_\sigma(a))$.
\end{definition}

\begin{proposition}[Order-independence of filtration]\label{prop:monoidal}
Let $\kappa = (d_1, \ldots, d_n)$ and let $\pi$ be any permutation
of $\{1, \ldots, n\}$. The predicate value $d_i(a)$ for any
element $a \in S$ depends only on $d_i$ and $a$, not on which
other discriminators have been evaluated or in what order.
The total profile of $a$ under $\kappa$ is the vector
$(d_1(a), \ldots, d_n(a)) \in \GF(2)^n$, independent of
evaluation order.
\end{proposition}

\begin{proof}
Each $d_i(a)$ is a function of $d_i$ and $a$ alone (by the
definition of a discriminator as a predicate $S \to \GF(2)$).
The profile vector is determined by $a$ and $\kappa$, not by
the sequence of evaluations.
\end{proof}

\begin{corollary}[Well-definedness of type erasure]
Projecting $\kappa$ to a sub-regime $\kappa' = (d_{i_1}, \ldots, d_{i_k})$
for $k < n$ yields a coarser but well-defined filtration.
\end{corollary}

\subsection{Membership as a relational property}

\begin{definition}[Membership]\label{def:membership}
An element $a$ is a \emph{member} of a collection $S$ under
a $\tau/\sigma$ pair $(d_\tau, d_\sigma)$ in regime $\kappa$,
written $a \in_{(d_\tau, d_\sigma),\kappa} S$,
if $d_\tau(a) = 1$ (the $\tau$-discriminator fires on $a$).
The element is in the \emph{complement} if $d_\sigma(a) = 1$.
Membership is a relation depending on the element, the
collection, and the discriminator pair.
\end{definition}

\begin{remark}[Yoneda content]
That membership depends on the collection --- not just the element
and the discriminator --- reflects the Yoneda perspective: the
identity of an element is constituted by its predicate profile
against the entire population under $\kappa$-filtration.
\end{remark}


\subsection{A concrete model}\label{sec:concrete-model}

We construct an explicit finite model to ground the
definitions. Let $V = \GF(2)^3$ and let the population
$S = \{a, b, c, d\}$ consist of four data.

\emph{Discriminators.}
Three basis discriminators $e_1, e_2, e_3 \in V$
(the standard basis of $\GF(2)^3$), with predicates:

\begin{center}
\small
\begin{tabular}{c|cccc|l}
& $a$ & $b$ & $c$ & $d$ & Keyed to \\
\hline
$e_1$ & 1 & 1 & 0 & 0 & $\tau$: ``has feature X'' \\
$e_2$ & 1 & 0 & 1 & 0 & $\sigma$: ``has feature Y'' \\
$e_3$ & 0 & 1 & 1 & 0 & $\tau$: ``has feature Z''
\end{tabular}
\end{center}

\emph{Profile linearity (Definition \ref{def:profile-linear}).} $e_1 + e_2$ is a new
discriminator in $V$. Its predicate:
$(e_1 + e_2)(a) = 1 + 1 = 0$,
$(e_1 + e_2)(b) = 1 + 0 = 1$,
$(e_1 + e_2)(c) = 0 + 1 = 1$,
$(e_1 + e_2)(d) = 0 + 0 = 0$.
This is the sum of the individual predicates, as required
by Definition \ref{def:profile-linear}. $\checkmark$

\emph{Membership profiles.}
Pairing $e_1$ ($\tau$-keyed) with $e_2$ ($\sigma$-keyed):
\begin{center}
\begin{tabular}{c|cc|l}
& $\tau = e_1$ & $\sigma = e_2$ & Profile \\
\hline
$a$ & 1 & 1 & $(1,1)$ overlap \\
$b$ & 1 & 0 & $(1,0)$ discriminated \\
$c$ & 0 & 1 & $(0,1)$ discriminated \\
$d$ & 0 & 0 & $(0,0)$ gap
\end{tabular}
\end{center}

All four cells are realized. Elements $b$ and $c$ are
discriminated (ordered). Element $a$ is in the overlap
(both features active); element $d$ is in the gap
(neither feature detected). The residue is
$R = \{a, d\}$: $\chi(a) = 1 + 1 = 0$,
$\chi(d) = 0 + 0 = 0$.

\emph{$\kappa$-evolution.}
The residue $\{a, d\}$ triggers evolution. Introducing
$e_3$ ($\tau$-keyed) and pairing with $e_2$ ($\sigma$-keyed)
on the residue: $e_3(a) = 0$, $e_2(a) = 1$, so $a$ gets
profile $(0,1)$ under the new pair --- discriminated.
$e_3(d) = 0$, $e_2(d) = 0$, so $d$ remains $(0,0)$ ---
still in the gap under this pair. A further evolution
would be needed for $d$.

\emph{The wedge product.}
$e_1 \wedge e_2 \neq 0$ (independent): the composition
of the $\tau$- and $\sigma$-keyed discriminators is
non-degenerate. $e_1 \wedge e_1 = 0$: self-composition
is degenerate. No discriminator can be paired with itself
to produce a non-trivial profile.

\emph{Russell exclusion.}
If $S \cap V \neq \varnothing$ (discriminators are also
data), the self-membership question is: does $e_1$,
viewed as a datum, carry the $\tau$ property that $e_1$
itself tests for? The Russell predicate would require
$\tau_{d_P}(e_i) = 1 + e_i(e_i)$ for each
discriminator-datum $e_i$. Since $\tau_{d_P}(0) = 1
\neq 0$ (the zero discriminator maps to $1$), this
profile is affine, violating profile linearity (Definition \ref{def:profile-linear}). $\checkmark$

This four-element model instantiates the full framework:
profile linearity, the four-cell structure, residue
detection, $\kappa$-evolution, the wedge product, and the
Russell exclusion. Every definition in \S\ref{sec:framework}
has a concrete realization.
